%============================================================
% CH1 · LESSON 2 — THE VOCABULARY OF GRAPHS
%============================================================
\subsection{The Vocabulary of Graphs}
\label{ssec:ch1-vocabulary}

Graph theory has a lot of words, but each one names a simple picture. This lesson is a picture dictionary: every term comes with a drawing, and the terms are ordered so that each one uses only the earlier ones. You do not have to memorise it in one go --- come back whenever a later chapter uses a word you are unsure of.

%------------------------------------------------------------
\subsubsection*{Degree}

The \textbf{degree} $\deg(v)$ of a vertex in an undirected graph is the number of edges touching it (a self-loop counts twice, because both of its ends touch $v$). In a simple graph this is the number of neighbours. A vertex of degree 0 is \textbf{isolated}; a vertex of degree 1 is a \textbf{leaf}.

We use one running example, the ``bow tie'' $G_1$ with $n=6$ and $m=7$:

\begin{diagrambox}
\begin{tikzpicture}[x=1.25cm, y=1cm]
    \node[gv] (1) at (0,0.8) {1};
    \node[gv] (2) at (0,-0.8) {2};
    \node[gvsel] (3) at (1.3,0) {3};
    \node[gvsel] (4) at (2.8,0) {4};
    \node[gv] (5) at (4.1,0.8) {5};
    \node[gv] (6) at (4.1,-0.8) {6};
    \foreach \a/\b in {1/2, 1/3, 2/3, 3/4, 4/5, 4/6, 5/6} \draw[ge] (\a) -- (\b);
    \foreach \v/\d/\pos in {1/2/left, 2/2/left, 3/3/above, 4/3/above, 5/2/right, 6/2/right}
        \node[gdist, \pos=1.5mm of \v] {deg $\d$};
\end{tikzpicture}
\figcap{$G_1$: edges 12, 13, 23, 34, 45, 46, 56. The two \legendchip{gvsel}{selected} vertices 3 and 4 have degree 3; all others have degree 2. Sum of degrees: $2+2+3+3+2+2=14=2\cdot 7$.}
\end{diagrambox}

In a \textbf{directed} graph each vertex has an \textbf{in-degree} $\indeg{v}$ (arcs entering $v$) and an \textbf{out-degree} $\outdeg{v}$ (arcs leaving $v$).

\begin{diagrambox}
\begin{tikzpicture}[x=1.4cm, y=1.1cm]
    \node[gv] (1) at (0,1) {1};
    \node[gv] (2) at (2,1) {2};
    \node[gv] (3) at (0,-0.6) {3};
    \node[gv] (4) at (2,-0.6) {4};
    \draw[gde] (1) -- (2);
    \draw[gde] (1) to[bend right=25] (3);
    \draw[gde] (3) to[bend right=25] (1);
    \draw[gde] (2) -- (3);
    \draw[gde] (3) -- (4);
    \draw[gde] (4) -- (2);
    \foreach \v/\i/\o/\pos in {1/1/2/left, 2/2/1/right, 3/2/2/left, 4/1/1/right}
        \node[gdist, \pos=1.5mm of \v, align=center] {in \i\\out \o};
    \node[note, text width=4.6cm, anchor=west] at (3.3,0.2) {arcs $1{\to}2$, $1{\to}3$, $2{\to}3$, $3{\to}1$, $3{\to}4$, $4{\to}2$\\[2pt]$\sum \text{in} = 1{+}2{+}2{+}1 = 6$\\$\sum \text{out} = 2{+}1{+}2{+}1 = 6 = m$};
\end{tikzpicture}
\figcap{Every arc has exactly one tail and one head, so the in-degrees and the out-degrees each add up to $m$.}
\end{diagrambox}

Degrees are the cheapest information a graph offers: one pass over the edge list computes all of them, without any search. Pattern~1 of this chapter solves whole problems with nothing else.

\begin{lstlisting}[style=cpp]
#include <bits/stdc++.h>
using namespace std;

// Degrees of an undirected graph + a check of the handshake lemma.
// Input: n m, then m lines "u v" (1-indexed).
int main() {
    int n, m;
    if (!(cin >> n >> m)) return 0;
    vector<int> deg(n + 1, 0);
    for (int i = 0; i < m; i++) {
        int u, v;
        cin >> u >> v;
        deg[u]++;            // the edge touches u ...
        deg[v]++;            // ... and v (a self-loop u==v adds 2 to u)
    }
    long long sum = 0;
    int odd = 0;
    for (int v = 1; v <= n; v++) {
        cout << "deg(" << v << ") = " << deg[v] << "\n";
        sum += deg[v];
        if (deg[v] % 2 == 1) odd++;
    }
    cout << "sum of degrees = " << sum << ", 2m = " << 2LL * m << "\n";
    cout << "odd-degree vertices = " << odd << "\n";
    return 0;
}
\end{lstlisting}

\begin{lstlisting}[style=trace]
input : 6 7   1 2  1 3  2 3  3 4  4 5  4 6  5 6        (the bow tie G1)
output: deg(1) = 2   deg(2) = 2   deg(3) = 3   deg(4) = 3   deg(5) = 2   deg(6) = 2
        sum of degrees = 14, 2m = 14
        odd-degree vertices = 2
\end{lstlisting}

%------------------------------------------------------------
\subsubsection*{The handshake lemma}

\begin{theorembox}[Handshake lemma]
In every undirected graph (self-loops and parallel edges allowed)
\[
    \sum_{v\in V} \deg(v) \;=\; 2\,|E|.
\]
\textbf{Proof.} Cut every edge in the middle into two \emph{half-edges}, one attached to each endpoint (a self-loop gives two half-edges at the same vertex). The degree of $v$ is exactly the number of half-edges at $v$, so the left side counts all half-edges. Every edge produced exactly two of them, so there are $2|E|$ half-edges. \hfill$\square$
\end{theorembox}

\begin{diagrambox}
\begin{tikzpicture}[x=1.25cm, y=1cm]
    \node[gv] (1) at (0,0.8) {1};
    \node[gv] (2) at (0,-0.8) {2};
    \node[gv] (3) at (1.3,0) {3};
    \node[gv] (4) at (2.8,0) {4};
    \node[gv] (5) at (4.1,0.8) {5};
    \node[gv] (6) at (4.1,-0.8) {6};
    \foreach \a/\b in {1/2, 1/3, 2/3, 3/4, 4/5, 4/6, 5/6}{
        \draw[draw=primaryblue, line width=1.4pt] (\a) -- ($(\a)!0.5!(\b)$);
        \draw[draw=accentorange, line width=1.4pt] ($(\a)!0.5!(\b)$) -- (\b);
        \fill[white] ($(\a)!0.5!(\b)$) circle (2.2pt);
    }
    \node[note, text width=4.2cm, anchor=west] at (5,0) {7 edges $\Rightarrow$ 14 half-edges.\\Each vertex owns as many half-edges as its degree.};
\end{tikzpicture}
\figcap{The proof in one picture: every edge of $G_1$ is cut into a blue and an orange half; counting halves by vertex gives $\sum\deg(v)$, counting them by edge gives $2m$.}
\end{diagrambox}

Three consequences used again and again:
\begin{enumerate}[leftmargin=*]
    \item \textbf{The number of odd-degree vertices is even.} The sum of all degrees is even; the even degrees contribute an even amount, so the odd degrees must add up to an even number, which needs an even count of them. ($G_1$ has two: vertices 3 and 4.)
    \item \textbf{Average degree $= 2m/n$.} A graph with $n=2\cdot10^5$ vertices and $m=2\cdot10^5$ edges has average degree 2 --- but a single vertex may still have degree $2\cdot10^5$ (a ``star''). Algorithms whose cost is $\sum_v \deg(v)$ are $\bigO{m}$; algorithms whose cost is $n\cdot\max\deg$ can be $\bigO{n^2}$.
    \item \textbf{Iterating over all adjacency lists costs $\bigO{n+m}$}, not $\bigO{n\cdot m}$: the lists have total length $\sum\deg(v)=2m$. This one fact is behind the running time of every traversal in this book.
\end{enumerate}

For directed graphs the same argument gives $\sum_v \indeg{v} = \sum_v \outdeg{v} = |E|$.

%------------------------------------------------------------
\subsubsection*{Walks, paths and cycles}

\begin{definitionbox}[Walk, path, cycle]
\begin{itemize}[leftmargin=*]
    \item A \textbf{walk} is a sequence of vertices $v_0, v_1, \dots, v_k$ in which consecutive vertices are joined by an edge (in a digraph: by an arc $v_{i}\to v_{i+1}$). Its \textbf{length} is the number of edges, $k$. Vertices and edges may repeat.
    \item A \textbf{path} (or \textbf{simple path}) is a walk with no repeated vertex.
    \item A \textbf{cycle} is a walk of length $\ge 3$ (in a digraph, $\ge 1$) that returns to its start, $v_k=v_0$, and repeats no other vertex.
\end{itemize}
\end{definitionbox}

\begin{diagrambox}
\begin{tikzpicture}[x=0.82cm, y=0.85cm]
    \foreach \sx/\title/\sub [count=\k] in {0/walk/{1, 2, 3, 1, 2}, 4.8/path/{1, 3, 4, 5}, 9.6/cycle/{4, 5, 6, 4}}{
        \begin{scope}[xshift=\sx cm]
            \node[gv] (a\k1) at (0,0.8) {1};
            \node[gv] (a\k2) at (0,-0.8) {2};
            \node[gv] (a\k3) at (1.3,0) {3};
            \node[gv] (a\k4) at (2.7,0) {4};
            \node[gv] (a\k5) at (4,0.8) {5};
            \node[gv] (a\k6) at (4,-0.8) {6};
            \foreach \a/\b in {1/2, 1/3, 2/3, 3/4, 4/5, 4/6, 5/6} \draw[gfaded] (a\k\a) -- (a\k\b);
            \node[font=\small\bfseries] at (2,-1.6) {\title};
            \node[font=\scriptsize\color{darkgray}] at (2,-2.05) {\sub};
        \end{scope}
    }
    \draw[gpath] (a11) -- (a12); \draw[gpath] (a12) -- (a13); \draw[gpath] (a13) -- (a11);
    \node[font=\scriptsize\color{accentorange}, anchor=west] at (-0.4,1.5) {edge 12 used twice};
    \draw[gpath] (a21) -- (a23); \draw[gpath] (a23) -- (a24); \draw[gpath] (a24) -- (a25);
    \draw[gpath] (a34) -- (a35); \draw[gpath] (a35) -- (a36); \draw[gpath] (a36) -- (a34);
\end{tikzpicture}
\figcap{On $G_1$: the walk 1, 2, 3, 1, 2 (length 4) repeats vertices; the path 1, 3, 4, 5 (length 3) does not; the cycle 4, 5, 6, 4 (length 3) returns to its start.}
\end{diagrambox}

\begin{insightbox}[Every walk contains a path]
If a walk from $u$ to $v$ visits some vertex twice, cut out the loop between the two visits; the result is a shorter walk from $u$ to $v$. Repeating this gives a path. So ``$v$ can be reached from $u$'' means the same whether we think of walks or paths --- and the \emph{shortest} walk is always a path.
\end{insightbox}

%------------------------------------------------------------
\subsubsection*{Connectivity and components}

An undirected graph is \textbf{connected} if every vertex can reach every other vertex. Otherwise it falls into \textbf{connected components}: maximal groups of vertices that can all reach each other. ``Maximal'' means you cannot add another vertex to the group and keep it connected. Every vertex belongs to exactly one component; an isolated vertex is a component on its own.

\begin{diagrambox}
\begin{tikzpicture}[x=1.1cm, y=1cm]
    \node[gvA] (1) at (0,0.7) {1};
    \node[gvA] (2) at (1.2,0.7) {2};
    \node[gvA] (3) at (0.6,-0.5) {3};
    \node[gvB] (4) at (2.8,0.7) {4};
    \node[gvB] (5) at (2.8,-0.6) {5};
    \node[gvsel] (6) at (4.4,0.1) {6};
    \node[gvseen] (7) at (5.9,0.7) {7};
    \node[gvseen] (8) at (7.1,0.7) {8};
    \node[gvseen] (9) at (6.5,-0.5) {9};
    \foreach \a/\b in {1/2, 2/3, 3/1, 4/5, 7/8, 8/9} \draw[ge] (\a) -- (\b);
    \foreach \x/\t in {0.6/{$\{1,2,3\}$}, 2.8/{$\{4,5\}$}, 4.4/{$\{6\}$}, 6.5/{$\{7,8,9\}$}}
        \node[font=\scriptsize] at (\x,-1.35) {\t};
\end{tikzpicture}
\figcap{A graph with $n=9$, $m=6$ and four components (one colour each). Vertex 6 is isolated: it is its own component. Counting and labelling components is Chapter~4; doing it with DSU is Chapter~6.}
\end{diagrambox}

For directed graphs ``connected'' splits into two notions: $v$ is \textbf{reachable} from $u$ if there is a directed path $u\leadsto v$, and the graph is \textbf{strongly connected} if every vertex is reachable from every other. A one-way road from A to B makes B reachable from A but not the other way round.

%------------------------------------------------------------
\subsubsection*{Trees and forests}

\begin{definitionbox}[Tree, forest]
A \textbf{tree} is a connected undirected graph with no cycle. A \textbf{forest} is a graph with no cycle (its components are trees). For a graph with $n$ vertices the following are equivalent --- any two of the first three imply the third:
\begin{itemize}[leftmargin=*]
    \item it is connected;\quad it has no cycle;\quad it has exactly $n-1$ edges;
    \item between any two vertices there is \textbf{exactly one} path.
\end{itemize}
A forest with $n$ vertices and $k$ trees has exactly $n-k$ edges.
\end{definitionbox}

\begin{diagrambox}
\begin{tikzpicture}[x=0.95cm, y=0.95cm]
    \begin{scope}
        \node[gv] (1) at (1,1) {1};
        \node[gv] (2) at (0,0) {2};
        \node[gv] (3) at (2,0) {3};
        \node[gv] (4) at (-0.5,-1.1) {4};
        \node[gv] (5) at (0.5,-1.1) {5};
        \node[gv] (6) at (2,-1.1) {6};
        \foreach \a/\b in {1/2, 1/3, 2/4, 2/5, 3/6} \draw[gtree] (\a) -- (\b);
        \node[font=\small\bfseries] at (0.75,-1.9) {tree};
        \node[font=\scriptsize\color{darkgray}] at (0.75,-2.3) {$n=6$, $m=5$};
    \end{scope}
    \begin{scope}[xshift=4.4cm]
        \node[gv] (1) at (1,1) {1};
        \node[gv] (2) at (0,0) {2};
        \node[gv] (3) at (2,0) {3};
        \node[gv] (4) at (-0.5,-1.1) {4};
        \node[gv] (5) at (0.5,-1.1) {5};
        \node[gv] (6) at (2,-1.1) {6};
        \foreach \a/\b in {1/2, 1/3, 2/4, 2/5, 3/6} \draw[gtree] (\a) -- (\b);
        \draw[gback, -] (5) -- (6);
        \node[font=\small\bfseries] at (0.75,-1.9) {add an edge};
        \node[font=\scriptsize\color{darkgray}] at (0.75,-2.3) {cycle 5--2--1--3--6--5};
    \end{scope}
    \begin{scope}[xshift=8.8cm]
        \node[gv] (1) at (1,1) {1};
        \node[gv] (2) at (0,0) {2};
        \node[gv] (3) at (2,0) {3};
        \node[gv] (4) at (-0.5,-1.1) {4};
        \node[gv] (5) at (0.5,-1.1) {5};
        \node[gv] (6) at (2,-1.1) {6};
        \foreach \a/\b in {1/2, 2/4, 2/5, 3/6} \draw[gtree] (\a) -- (\b);
        \draw[gcut] (1.1,0.2) -- (1.9,0.9);
        \node[font=\small\bfseries] at (0.75,-1.9) {remove an edge};
        \node[font=\scriptsize\color{darkgray}] at (0.75,-2.3) {forest of 2 trees, $m=4$};
    \end{scope}
\end{tikzpicture}
\figcap{A tree is ``just connected'': adding any edge closes exactly one cycle (dashed), removing any edge disconnects it. That is why $n-1$ edges is the magic number.}
\end{diagrambox}

When one vertex is chosen as the \textbf{root}, every other vertex gets a \textbf{parent} (its neighbour towards the root) and \textbf{children}; the tree becomes a \textbf{rooted tree}. Company hierarchies, file systems and DFS trees are rooted trees. Trees are so important that Chapter~13 is devoted to them.

%------------------------------------------------------------
\subsubsection*{DAGs}

A directed graph with no directed cycle is a \textbf{DAG} (directed acyclic graph). Prerequisites, task dependencies and ``$i<j$'' relations are DAGs. Their vertices can be listed so that every arc points forward --- a \textbf{topological order} (Chapter~5) --- and that is exactly the order in which DP fills its table.

\begin{diagrambox}
\begin{tikzpicture}[x=1.15cm, y=1cm]
    \begin{scope}
        \node[gv] (a) at (0,0.8) {a};
        \node[gv] (b) at (1.5,0.8) {b};
        \node[gv] (c) at (0,-0.6) {c};
        \node[gv] (d) at (1.5,-0.6) {d};
        \node[gv] (e) at (3,0.1) {e};
        \draw[gde] (a) -- (b); \draw[gde] (a) -- (c); \draw[gde] (c) -- (d);
        \draw[gde] (b) -- (d); \draw[gde] (b) -- (e); \draw[gde] (d) -- (e);
        \node[font=\small\bfseries] at (1.5,-1.4) {a DAG};
    \end{scope}
    \begin{scope}[xshift=4.8cm]
        \foreach \v/\x in {a/0, c/1.1, b/2.2, d/3.3, e/4.4} \node[gv] (\v) at (\x,0.1) {\v};
        \draw[gde] (a) -- (c); \draw[gde] (b) -- (d); \draw[gde] (d) -- (e);
        \draw[gde] (a) to[bend left=35] (b);
        \draw[gde] (b) to[bend left=35] (e);
        \draw[gde] (c) to[bend right=35] (d);
        \node[font=\small\bfseries] at (2.2,-1.4) {same DAG, topological order};
        \node[font=\scriptsize\color{darkgray}] at (2.2,-1.8) {every arc points to the right};
    \end{scope}
    \begin{scope}[xshift=11.5cm]
        \node[gv] (x) at (0,0.8) {x};
        \node[gv] (y) at (1.4,0.8) {y};
        \node[gv] (z) at (0.7,-0.5) {z};
        \draw[gde] (x) -- (y); \draw[gde] (y) -- (z); \draw[gback] (z) -- (x);
        \node[font=\small\bfseries] at (0.7,-1.4) {not a DAG};
    \end{scope}
\end{tikzpicture}
\figcap{A DAG can be laid out on a line with all arcs going right. The cycle $x\to y\to z\to x$ makes such a layout impossible: whichever of the three comes first has an arc coming \emph{into} it from the right.}
\end{diagrambox}

%------------------------------------------------------------
\subsubsection*{Bipartite graphs}

A graph is \textbf{bipartite} if its vertices can be coloured with two colours so that every edge joins two \emph{different} colours --- for example, students and the clubs they belong to, or the black and white squares of a chessboard (a knight always changes colour). A graph is bipartite exactly when it has \textbf{no cycle of odd length} (Chapter~4).

\begin{diagrambox}
\begin{tikzpicture}[x=1.1cm, y=1cm]
    \begin{scope}
        \node[gvA] (1) at (0,1) {1};
        \node[gvB] (2) at (1.3,1) {2};
        \node[gvA] (3) at (1.3,-0.3) {3};
        \node[gvB] (4) at (0,-0.3) {4};
        \node[gvB] (5) at (2.6,-0.3) {5};
        \foreach \a/\b in {1/2, 2/3, 3/4, 4/1, 3/5} \draw[ge] (\a) -- (\b);
        \node[font=\small\bfseries] at (1.3,-1.2) {bipartite};
    \end{scope}
    \begin{scope}[xshift=4.5cm]
        \node[gvA] (1) at (0,1) {1};
        \node[gvA] (3) at (0,0.1) {3};
        \node[gvB] (2) at (2,1.2) {2};
        \node[gvB] (4) at (2,0.35) {4};
        \node[gvB] (5) at (2,-0.5) {5};
        \foreach \a/\b in {1/2, 3/2, 3/4, 1/4, 3/5} \draw[ge] (\a) -- (\b);
        \node[font=\small\bfseries] at (1,-1.2) {same graph, two sides};
    \end{scope}
    \begin{scope}[xshift=9.2cm]
        \node[gvA] (a) at (0,1) {a};
        \node[gvB] (b) at (1.4,1) {b};
        \node[gvdead] (c) at (0.7,-0.2) {c};
        \draw[ge] (a) -- (b); \draw[ge] (b) -- (c); \draw[gcut] (c) -- (a);
        \node[font=\scriptsize\color{warningred}, align=center] at (2.5,-0.2) {c needs a colour\\different from both};
        \node[font=\small\bfseries] at (0.7,-1.2) {odd cycle: not bipartite};
    \end{scope}
\end{tikzpicture}
\figcap{Left and middle: a bipartite graph redrawn with its two sides as columns --- every edge crosses. Right: a triangle cannot be 2-coloured.}
\end{diagrambox}

%------------------------------------------------------------
\subsubsection*{Complete graphs and complements}

The \textbf{complete graph} $K_n$ has an edge between \emph{every} pair of its $n$ vertices, so $m=\binom{n}{2}=\frac{n(n-1)}{2}$, the maximum for a simple undirected graph. The \textbf{complement} $\overline{G}$ of a simple graph $G$ has the same vertices and exactly the edges that $G$ does \emph{not} have; together, $G$ and $\overline{G}$ form $K_n$.

\begin{diagrambox}
\begin{tikzpicture}[x=1cm, y=1cm]
    \foreach \k/\sx in {3/0, 4/2.5, 5/5.3}{
        \begin{scope}[xshift=\sx cm]
            \foreach \i in {1,...,\k} \node[gv, minimum size=5.5mm, font=\scriptsize\bfseries] (k\k-\i) at ({90+360/\k*(\i-1)}:0.9) {\i};
            \foreach \i in {1,...,\k} \foreach \j in {1,...,\k} {
                \ifnum\i<\j \draw[ge] (k\k-\i) -- (k\k-\j); \fi }
            \pgfmathtruncatemacro{\mm}{\k*(\k-1)/2}
            \node[font=\scriptsize\bfseries] at (0,-1.45) {$K_{\k}$: $m=\mm$};
        \end{scope}
    }
    \begin{scope}[xshift=8.4cm, yshift=0.3cm]
        \foreach \i/\x in {1/0, 2/0.9, 3/1.8, 4/2.7} \node[gv, minimum size=5.5mm, font=\scriptsize\bfseries] (p\i) at (\x,0.5) {\i};
        \draw[ge] (p1) -- (p2); \draw[ge] (p2) -- (p3); \draw[ge] (p3) -- (p4);
        \foreach \i/\x in {3/0, 1/0.9, 4/1.8, 2/2.7} \node[gv, minimum size=5.5mm, font=\scriptsize\bfseries] (q\i) at (\x,-0.9) {\i};
        \draw[gpath, line width=1.2pt] (q3) -- (q1); \draw[gpath, line width=1.2pt] (q1) -- (q4); \draw[gpath, line width=1.2pt] (q4) -- (q2);
        \node[font=\scriptsize, anchor=west] at (3.0,0.5) {$G$};
        \node[font=\scriptsize, anchor=west] at (3.0,-0.9) {$\overline{G}$};
    \end{scope}
\end{tikzpicture}
\figcap{$K_3$, $K_4$, $K_5$ with 3, 6 and 10 edges. Right: the path $G=1$--$2$--$3$--$4$ (3 edges) and its complement (the other $6-3=3$ pairs: 13, 14, 24), redrawn: it is again a path, 3--1--4--2.}
\end{diagrambox}

Complements matter in contests because a statement may describe the edges a graph does \emph{not} have (``every two cities are connected except for these $m$ pairs''). If $n$ is large, the complement is dense and cannot be built explicitly --- a classic trap.

%------------------------------------------------------------
\subsubsection*{Sparse and dense graphs}

A simple graph has between $0$ and $\binom n2\approx n^2/2$ edges. When $m$ is of the order of $n$ (trees, road maps, grids, typical input with $n,m\le 2\cdot10^5$) the graph is \textbf{sparse}; when $m$ is of the order of $n^2$ it is \textbf{dense}. The distinction decides how to store the graph (Lesson~\ref{ssec:ch1-storing}) and which algorithm is fastest.

\begin{table}[H]
\centering\small
\begin{tabular}{@{}l l l l@{}}
\toprule
\textbf{Graph} & \textbf{$m$ as a function of $n$} & \textbf{$n=1000$} & \textbf{kind} \\
\midrule
tree, path, functional graph & $n-1$ or $n$ & $\approx 10^3$ & sparse \\
grid $r\times c$ (4 neighbours) & $\approx 2n$ & $\approx 2\cdot 10^3$ & sparse \\
typical contest input & given, usually $\le 2\cdot 10^5$ & --- & sparse \\
complete graph $K_n$ & $n(n-1)/2$ & $499\,500$ & dense \\
complete digraph & $n(n-1)$ & $999\,000$ & dense \\
\bottomrule
\end{tabular}
\end{table}

%------------------------------------------------------------
\subsubsection*{Functional graphs: exactly one arrow out}

A directed graph in which \textbf{every vertex has out-degree exactly 1} is a \textbf{functional graph}: it is the picture of a function $f:\{1..n\}\to\{1..n\}$, with an arc $v\to f(v)$. Permutations, ``each student blames one other student'', ``teleporter $i$ sends you to $a_i$'' are all functional graphs. It has exactly $n$ arcs, and its shape is always the same: \emph{every component is one cycle with trees hanging into it}.

\begin{diagrambox}
\begin{tikzpicture}[x=1.15cm, y=1cm]
    \node[gvsel] (2) at (2,0.8) {2};
    \node[gvsel] (3) at (3.2,0) {3};
    \node[gvsel] (4) at (2,-0.8) {4};
    \node[gv] (1) at (0.8,1.5) {1};
    \node[gv] (5) at (0.8,-0.8) {5};
    \node[gv] (6) at (-0.4,-0.8) {6};
    \node[gvsrc] (7) at (-1.6,-0.8) {7};
    \draw[gde] (2) to[bend left=20] (3);
    \draw[gde] (3) to[bend left=20] (4);
    \draw[gde] (4) to[bend left=20] (2);
    \draw[gde] (1) -- (2); \draw[gde] (5) -- (4); \draw[gde] (6) -- (5); \draw[gde] (7) -- (6);
    \node[gvsel] (8) at (5,0.4) {8};
    \node[gvsel] (9) at (6.4,0.4) {9};
    \draw[gde] (8) to[bend left=30] (9);
    \draw[gde] (9) to[bend left=30] (8);
    \node[note, text width=4cm] at (5.7,-0.9) {$f = (2,3,4,2,4,5,6,9,8)$\\for $v = 1..9$};
\end{tikzpicture}
\figcap{$f(1)=2, f(2)=3, f(3)=4, f(4)=2, f(5)=4, f(6)=5, f(7)=6, f(8)=9, f(9)=8$. Two components; the \legendchip{gvsel}{cycle} vertices are $\{2,3,4\}$ and $\{8,9\}$. Starting at 7 and following arrows: 7, 6, 5, 4, 2, 3, 4, 2, 3, 4, \ldots{} --- a tail of length 3, then the cycle forever (the shape of the Greek letter $\rho$).}
\end{diagrambox}

\begin{insightbox}[Why there is always a cycle]
Follow the arrows from any vertex. There are only $n$ vertices, so within $n$ steps some vertex repeats; from then on the walk repeats forever. Pattern~5 turns this observation into algorithms, and Chapter~11 (functional graphs and binary lifting) jumps $10^{18}$ steps along such walks.
\end{insightbox}

%------------------------------------------------------------
\subsubsection*{The dictionary at a glance}

\begin{table}[H]
\centering\small
\begin{tabularx}{\linewidth}{@{}l X l@{}}
\toprule
\textbf{Term} & \textbf{Meaning} & \textbf{Typical clue in a statement} \\
\midrule
degree & number of incident edges & ``how many friends'', ``number of roads at'' \\
path / cycle & no repeated vertex / returns to start & ``route'', ``loop'', ``come back'' \\
component & maximal mutually reachable group & ``groups'', ``islands'', ``clusters'' \\
tree & connected, $n-1$ edges, no cycle & ``$n-1$ roads'', ``unique path'' \\
DAG & directed, no directed cycle & ``prerequisites'', ``before'', ``$i<j$'' \\
bipartite & 2-colourable, no odd cycle & ``two teams'', ``enemies in different groups'' \\
complete / complement & all pairs / the missing pairs & ``every pair except'' \\
functional graph & out-degree exactly 1 & ``$a_i$'', ``next'', ``teleports to'' \\
\bottomrule
\end{tabularx}
\end{table}

%------------------------------------------------------------
\subsubsection*{Check your understanding}

\begin{enumerate}[leftmargin=*]
    \item Is there a simple graph with degrees $3,3,3,2,2$?
    \item Nine people at a party: can each shake hands with exactly three others?
    \item A connected graph has $n=10$ vertices and $m=9$ edges. What is it? What if it had $m=10$ edges?
    \item A forest has 20 vertices and 16 edges. How many trees does it have?
    \item Is the cycle on 6 vertices bipartite? The cycle on 7 vertices?
    \item How many edges does the complement of a tree on $n$ vertices have?
    \item A functional graph has $n$ vertices. How many arcs does it have, and can it be a DAG?
\end{enumerate}

\begin{solutionbox}[Answers]
\begin{enumerate}[leftmargin=*]
    \item No: the degrees add up to 13, which is odd, but $\sum\deg = 2m$ is even (equivalently: three odd degrees, an odd number).
    \item No, for the same reason: $9\cdot 3=27$ handshake-ends cannot be split into pairs.
    \item With 9 edges it is a tree. With 10 edges it has exactly one cycle (a tree plus one edge; sometimes called a \emph{unicyclic} graph).
    \item $n-k=m$ gives $k=20-16=4$ trees.
    \item $C_6$ is bipartite (colour alternately); $C_7$ is not --- it is an odd cycle.
    \item $\binom n2-(n-1)=\frac{(n-1)(n-2)}{2}$.
    \item Exactly $n$ arcs (one per vertex). It is never a DAG: following arrows must eventually repeat a vertex, so every functional graph contains a cycle.
\end{enumerate}
\end{solutionbox}
